\documentclass[12pt]{article}
\usepackage[utf8]{vietnam}
\usepackage{amsmath, amssymb}
\usepackage{enumitem}

\begin{document}

\section*{Bài tập 10}
Đường thẳng $d$ có một vectơ pháp tuyến là $\vec{n} = (-2; -5)$. Đường thẳng $\Delta$ song song với $d$. Tìm một vectơ chỉ phương của đường thẳng $d$.

\section*{BÀI TẬP TRẮC NGHIỆM}
\subsection*{PHẦN I. Câu trắc nghiệm nhiều phương án lựa chọn. Mỗi câu hỏi thí sinh chỉ chọn một phương án.}

\begin{enumerate}[label=\textbf{Câu \arabic*:}]
\item Trong mặt phẳng Oxy, cho đường thẳng $d: 2x - 3y + 1 = 0$. Một vectơ pháp tuyến của đường thẳng $d$ là:
\[
\begin{array}{ll}
\text{A. } \vec{n} = (2; -3) & \text{B. } \vec{n} = (3; 2) \\
\text{C. } \vec{n} = (3; -2) & \text{D. } \vec{n} = (2; 3)
\end{array}
\]

\item Trong mặt phẳng với hệ toạ độ Oxy, một vectơ chỉ phương của đường thẳng $d: \begin{cases} x = 1 - 2t \\ y = 2 + 3t \end{cases}$ là:
\[
\begin{array}{ll}
\text{A. } \vec{a} = (2; 3) & \text{B. } \vec{b} = (3; 2) \\
\text{C. } \vec{c} = (3; -2) & \text{D. } \vec{d} = (2; -3)
\end{array}
\]

\item Trong mặt phẳng Oxy, cho đường thẳng $(d)$ có phương trình $2x - y + 5 = 0$. Tìm một vectơ chỉ phương của $(d)$:
\[
\text{A. } (1; -2) \qquad \text{B. } (2; 1) \qquad \text{C. } (2; -1) \qquad \text{D. } (1; 2)
\]

\item Trong mặt phẳng Oxy, cho đường thẳng $d: 3x - 7y - 1 = 0$. Vectơ nào sau đây là vectơ pháp tuyến của đường thẳng $d$:
\[
\text{A. } \vec{n} = (3; -7) \qquad \text{B. } \vec{n} = (2; 3) \qquad \text{C. } \vec{n} = (3; 7) \qquad \text{D. } \vec{n} = (7; 3)
\]

\item Trong mặt phẳng Oxy, cho đường thẳng $d: y = -3x + 5$. Một vectơ pháp tuyến của đường thẳng $d$ là:
\[
\text{A. } \vec{n} = (1; 3) \qquad \text{B. } \vec{n} = (3; 1) \qquad \text{C. } \vec{n} = (-3; 1) \qquad \text{D. } \vec{n} = (1; -3)
\]

\item Trong mặt phẳng Oxy, cho đường thẳng $d: 4x + 5y - 4 = 0$. Vectơ nào sau đây \textbf{không phải} là vectơ pháp tuyến của đường thẳng $d$:
\[
\text{A. } \vec{n}_1 = (4; 5) \qquad \text{B. } \vec{n}_2 = (-8; -10) \qquad \text{C. } \vec{n}_3 = (4; -5) \qquad \text{D. } \vec{n}_4 = \left(\frac{4}{3}; \frac{5}{3}\right)
\]

\item Trong mặt phẳng Oxy, cho hai điểm $A(1;1)$, $B(2;3)$. Tìm một vectơ pháp tuyến của đường trung trực của đoạn thẳng $AB$:
\[
\text{A. } \vec{n} = (1; -2) \qquad \text{B. } \vec{n} = (2; 1) \qquad \text{C. } \vec{n} = (-1; 2) \qquad \text{D. } \vec{n} = (1; 2)
\]

\item Trong mặt phẳng Oxy, cho tam giác $ABC$ biết $A(-1;-1)$, $B(1;-3)$, $C(2;4)$. Tìm một vectơ pháp tuyến của đường cao kẻ từ $B$ của tam giác $ABC$:
\[
\text{A. } \vec{n} = (3; 5) \qquad \text{B. } \vec{n} = (3; -5) \qquad \text{C. } \vec{n} = (5; 3) \qquad \text{D. } \vec{n} = (-5; 3)
\]

\item Trong mặt phẳng Oxy, vectơ nào dưới đây là một vectơ pháp tuyến của $d: x - 2y + 2025 = 0$?
\[
\text{A. } \vec{n}_1 = (0; -2) \qquad \text{B. } \vec{n}_3 = (-2; 0) \qquad \text{C. } \vec{n}_4 = (2; 1) \qquad \text{D. } \vec{n}_2 = (1; -2)
\]

\item Cho đường thẳng $(d): \begin{cases} x = 2 - 3t \\ y = -1 + 2t \end{cases}$ và điểm $A\left(\dfrac{7}{2}; -2\right)$. Điểm $A \in (d)$ ứng với giá trị nào của $t$?
% (Câu hỏi này trong file gốc không có các phương án lựa chọn, giữ nguyên đề)
\end{enumerate}

\end{document}